\section{Discussion}
\label{sec:conclusion}

\paragraph{Lessons for quantum components in continual learning.}
\begin{enumerate}
  \item \emph{Measure the noise first.} The baseline moves by $0.14$--$0.38$
    points between runs. An earlier $+0.36$ gain of ours disappeared once
    measured properly. Evaluating all inference heads on the same network
    reduces the noise of a comparison to $0.03$--$0.05$ points; it is cheap,
    and it should be standard whenever a component changes only inference.
  \item \emph{Always test the classical twin.} Beating the baseline was easy:
    the PGM fusion passed that test. Beating the classical method that sees
    the same inputs was not. Without the twin, Table~\ref{tab:main} would
    have read as a quantum improvement.
  \item \emph{Know which degrees of freedom are real.} For a measurement on
    amplitude-encoded states, trainable circuits after the encoding and
    state shrinkage do nothing (Propositions~\ref{prop:unitary}
    and~\ref{prop:shrink}). Only the encoding remains. Effort spent on
    deeper circuits in such a pipeline cannot pay off.
  \item \emph{Local kernels are fragile.} Entangling layers spread the
    information that the projected kernel needs into correlations it does
    not read, which is a concrete form of concentration.
\end{enumerate}

\paragraph{What did help.}
The state-discrimination view led to a useful classical result. The
statistics kept by exemplar-free methods carry information that the aligned
classifier does not use, and a calibrated fusion recovers part of it.
Final accuracy rises on all four benchmarks, without training or new memory.
The gain is bounded: about $3.4$ points of the gap to joint training lie in
the representation, which no inference-time rule can change.

\paragraph{Limitations.}
\begin{itemize}
  \item We study one baseline and four benchmarks with three seeds each.
    The transfer test failed on CUB-200.
  \item Our circuits are small and simulated. The negative results concern
    these placements and encodings, not quantum models in general.
  \item A quantum component that changes the representation, such as a
    measurement-based training loss, remains untested. It is the natural
    next step, again with its classical twin fixed in advance.
\end{itemize}

\paragraph{Reproducibility.}
Configurations, logs, analysis scripts and the pre-registered
specification of every round are in the project repository under
\texttt{exps/}. Every number in this paper is generated from these files.
